\section*{Chapter \ref{ch:ml:targets-labels-and-sample-weights} --- Targets, Labels and Sample Weights}

\begin{solution}{exo:ml:targets-labels-and-sample-weights:1}
About $1\,000/5 = 200$ independent labels (\cref{prop:ml:labels:sum}); uniqueness $1/5 = 0.2$ away from the ends.
\end{solution}

\begin{solution}{exo:ml:targets-labels-and-sample-weights:2}
Bars 1 to 8: $c = 1, 2, 2, 2, 0, 0, 1, 1$. Uniqueness: $(1 + \tfrac12)/2 = 0.75$; $(\tfrac12 + \tfrac12 + \tfrac12)/3 =
0.5$; $(\tfrac12 + \tfrac12)/2 = 0.5$; $1$.
\end{solution}

\begin{solution}{exo:ml:targets-labels-and-sample-weights:3}
$1.5\times1.2\%\times\sqrt{16} = 7.2\%$ each side.
\end{solution}

\begin{solution}{exo:ml:targets-labels-and-sample-weights:4}
Trades opened daily with 20-day spans have uniqueness about $1/20$, so the effective count is about 1\,250 and the
$t$-statistic about $6.1/\sqrt{20} = 1.36$: not significant.
\end{solution}

\begin{solution}{exo:ml:targets-labels-and-sample-weights:5}
The $t$-statistic is the mean over the standard deviation times the square root of the effective number of trades.
Keeping a third of the trades raises the mean by half (0.064 to 0.097 widths) but divides the effective count by about
three, so $t$ falls (3.48 to 3.19). It would rise if the secondary model found trades with negative expectation to drop
(information the primary model lacks), not merely trades with a smaller positive one.
\end{solution}

\begin{solution}{exo:ml:targets-labels-and-sample-weights:6}
The whole-sample standard deviation uses the future (look-ahead) and ignores volatility clustering: barriers are too
wide in calm periods and too narrow in stress, so labels settle for reasons unrelated to the trade. Use the volatility
known at the event (EWMA or GARCH, Book~4, chapter~18).
\end{solution}

\begin{solution}{exo:ml:targets-labels-and-sample-weights:7}
\texttt{label\_stats(k=2.0)}: profit 4.0\%, stop 3.3\%, vertical 92.7\%; labels last 9.8 days, uniqueness 0.1025
(0.123 at $k = 1$). Wider barriers are rarely touched within ten days, so the barrier label converges to the
\omterm{def:ml:targets-labels-and-sample-weights:fixed}{fixed-horizon label}.
\end{solution}

\begin{solution}{exo:ml:targets-labels-and-sample-weights:8}
With constant uniqueness, cumulative uniqueness is proportional to the label's rank: the halfway label has weight
$0.5 + 0.5\times0.5 = 0.75$, and the average of a linear ramp from 0.5 to 1 is 0.75. With varying uniqueness, crowded
periods (low uniqueness) advance the ramp slowly, so their many labels share similar weights and the decay follows
information, not the calendar; the mean then depends on where the crowded periods fall.
\end{solution}

\begin{solution}{pb:ml:targets-labels-and-sample-weights:1}
\begin{enumerate}
\item Profit 22.7\%, stop 20.3\%, vertical 57.0\%.
\item 8.3 days: a label settles at the first touch, before day ten in 43\% of the events.
\item Fixed horizon: up (a winner, $+1.84\%$); barrier: stopped at $-5.04\%$ on day three, a loser. The barrier label
describes the trade.
\item 1.9\%; sizes differ for every touched label, since barrier returns are truncated as the trade's are.
\item Concurrency 10, uniqueness 0.100.
\item $20\times2\,458/10 = 4\,916$; barrier labels 6\,040.
\item 9.4 on rows, 3.3 adjusted; ratio $2.85 = \sqrt{1/0.123}$ (\cref{prop:ml:labels:sum}, part 2).
\item Standard errors, $t$-statistics, bootstrap intervals, out-of-bag scores, and any significance test that counts
rows.
\item 0.65 and 0.69.
\item Plain 0.504, 0.517, 0.509; small 0.511, 0.512, 0.506; weighted 0.506, 0.511, 0.510.
\item The trees were regularised and no feature identifies the date, so near-duplicates in a bootstrap sample could not
be memorised into an advantage.
\item Deep trees with a slowly moving feature that locates time (a regime variable, a trend of the level): overlapping
neighbours then share both features and labels, and full-size bootstrap samples fit them.
\item Hit rate 54.5\% against 53.0\% and 0.097 against 0.064 widths per trade, on a third of the trades; the
overlap-adjusted $t$ falls from 3.48 to 3.19.
\item It chooses which of the primary model's trades to take and how much; the side, where the signal is, comes from
the primary model.
\item Yes, in proportion to $2p - 1$ or to a Kelly-like fraction (chapter~11), which keeps every trade but scales it:
filtering is the special case of sizes 0 and 1.
\item \emph{Named result}: 4\,916 effective labels of 48\,980 (10\%); $t$ of 9.4 on rows, 3.3 adjusted; uniqueness
weighting's AUC gain over plain bagging $-0.001$ on average, within $\pm0.007$ market by market.
\item About $20\times2\,450/20 = 2\,450$: half as many.
\item Rows, spans ($t_0$, $t_1$), the sum of uniqueness, and the number of independent dates.
\item Both come from the spans: purging removes training labels whose spans overlap the test fold; uniqueness measures
how much spans overlap inside a sample.
\item A label is the settlement of a trade: what the position made, from the decision time to the exit.
\end{enumerate}
\end{solution}

\begin{solution}{iq:ml:targets-labels-and-sample-weights:1}
It ignores the path (stops and targets), ignores volatility (a fixed return threshold means different things in
different regimes), and consecutive labels overlap in 19 of 20 days, so the sample is worth about a twentieth of its
rows.
\iqlookfor{path, volatility and overlap.}
\end{solution}

\begin{solution}{iq:ml:targets-labels-and-sample-weights:2}
For each event, set a profit barrier and a stop at multiples of the current volatility, and a vertical barrier at a
maximum holding time; the label is the outcome at the first barrier touched, and the label records when it settled.
\iqlookfor{volatility scaling, the first touch, and the settlement time kept for purging.}
\end{solution}

\begin{solution}{iq:ml:targets-labels-and-sample-weights:3}
Per stock, about $2\,520/10 = 252$ independent labels, so about 100\,000 in total, and fewer since stocks are correlated
(a common factor makes a date worth a few independent stocks). The effective sample favours regularised, low-variance
models and honest standard errors.
\iqlookfor{uniqueness times rows, then cross-sectional correlation.}
\end{solution}

\begin{solution}{iq:ml:targets-labels-and-sample-weights:4}
A secondary classifier estimates the probability that a primary model's call is right, trained on labels of whether
each call made money; its output filters or sizes the trades. Useful when the primary side comes from a source the
model cannot reproduce (a discretionary view, a rule) and the secondary model sees when it works.
\iqlookfor{side versus size, and that it cannot add a signal the primary model lacks.}
\end{solution}

\begin{solution}{iq:ml:targets-labels-and-sample-weights:5}
Out-of-bag observations are the ones a tree did not draw, but with overlapping labels their neighbours (sharing most of
the label and similar features) were drawn: the score is not out of sample. The later test period has no such
neighbours. Use purged, time-ordered validation instead.
\iqlookfor{the leak from overlap into the out-of-bag estimate.}
\end{solution}

\begin{solution}{iq:ml:targets-labels-and-sample-weights:6}
Each innovation enters $h$ of the $n$ sums, so the mean of the sums is about $\frac hn\sum_t\varepsilon_t$, of variance
$h^2\sigma^2/n$; treating the sums as independent gives $\Var(h\text{-sum})/n = h\sigma^2/n$. The ratio is $h$.
\iqlookfor{the counting argument and its consequence for $t$-statistics ($\sqrt h$).}
\end{solution}
